\section{Data and Training Pipeline}

\subsection{Base Data}

The final active training regime uses Oval-only FATROP-clean data. The documented dataset status is:
\begin{itemize}
\item shift shard: 31,710 episodes,
\item hard-repair shard: 400 accepted episodes,
\item post-projection repair shard: 1,000 accepted episodes,
\item failure-conditioned repair shard: 300 accepted episodes.
\end{itemize}

The shift shard provides broad nominal coverage; the repair shards target exactly the off-manifold states that appear during DT rollout and wrapper correction.

\subsection{Training Runs}

Table~\ref{tab:training-runs} summarizes the main completed runs from the final repo state.

\begin{table}[t]
\centering
\caption{Main DT runs in the final repo state. Validation numbers are best logged checkpoints from \texttt{metrics.jsonl}.}
\label{tab:training-runs}
\begin{tabular}{lcccc}
\toprule
Run & Data Mix & $\lambda_x$ & Best Val Loss & Best Val Action Loss \\
\midrule
Parent improved & shift 99.5\%, hard 0.5\% & 0.10 & 0.002835 & 0.002600 \\
Post-proj FT2 rerun & shift 80\%, hard 5\%, postproj 15\% & 0.10 & 0.003170 & 0.002868 \\
Lateral FT & shift 75\%, hard 5\%, postproj 20\% & 0.05 & 0.002773 & 0.002544 \\
Failure-conditioned FT & shift 70\%, failure 20\%, hard 5\%, postproj 5\% & 0.05 & 0.002364 & 0.002197 \\
\bottomrule
\end{tabular}
\end{table}

The parent improved run used a larger architecture than earlier baselines: context length 50, 6 layers, 4 heads, and embedding size 192. Fine-tune runs reused that architecture and resumed from the parent checkpoint.

Two observations matter. First, offline validation clearly improved over the course of the project, especially after failure-conditioned data was introduced. Second, these offline improvements did not reliably predict downstream warm-start quality.

\subsection{Failure-Conditioned Repairs}

The failure-conditioned pipeline is one of the most substantive repo contributions. The process is:
\begin{enumerate}
\item export rollout traces from current checkpoints,
\item identify failure buckets such as repeated \texttt{e\_clip} and early lateral instability,
\item rebuild repair segments around those failures using the nonlinear solver, and
\item fine-tune on the resulting accepted segments.
\end{enumerate}

The documented trace export produced at least two populated buckets:
\begin{itemize}
\item \texttt{repeated\_e\_clip}: 2,987 examples,
\item \texttt{early\_lateral\_failure}: 162 examples.
\end{itemize}

The first completed failure-conditioned shard contained 300 accepted repair segments.
